\chapter{Cadeias de Prompts, Agentes e Workflows}
\label{ch:agentes}

\ficcao O assistente da Lavandaria recebe «Quando é que me devolvem a meia?». Pode procurar a política, pedir o talão ou encaminhar. Se a regra já diz que falta o talão, tens uma escolha de desenho. Deixas o modelo decidir o próximo passo ou o próprio programa segue a regra?

\section{O caminho que já conheces}

Uma \emph{cadeia de prompts} liga chamadas sucessivas. Para preparar uma página de apresentação, podes clarificar público e objetivo, escrever uma estrutura, preparar o texto e revê-lo. Cada etapa deve entregar algo que a seguinte consiga usar. Por exemplo, o rascunho recebe os factos e a estrutura. Se receber só uma síntese, uma condição perdida na síntese já não pode ser recuperada por escrever melhor.

Um \emph{workflow}, ou fluxo de trabalho, liga etapas com regras de passagem. Pode ser linear, escolher um ramo ou repetir um passo. Num atendimento, o programa pode consultar um registo para pedidos de estado e escolher outro ramo quando precisa de uma política. Uma validação pode interromper o fluxo antes do envio. Um ciclo precisa de uma condição de saída e de um orçamento de tentativas. Uma nova tentativa faz sentido quando há uma correção a fazer ou um erro transitório identificado.

\begin{promptlabbox}{Separar um pedido grande em etapas verificáveis}
\textbf{Objetivo e contexto:} Comparar um pedido único e um fluxo com etapas. Usa o mesmo assunto, fontes, público e critérios.

Prepara um texto a partir de uma nota factual. Numa versão, pede o texto de uma vez. Na outra, obtém uma estrutura, revê as condições e só depois pede o rascunho. Faz as duas versões manualmente ou em sessões separadas. Guarda as entradas e o resultado de cada etapa.

\textbf{Resultado esperado e validação:} As duas versões são avaliadas pela mesma regra. Localiza em que etapa se perdeu ou se preservou cada condição. Regista tempo e chamadas se usaste um modelo.

\textbf{Falhas possíveis:} Dar mais factos só a uma das versões, ocultar falhas intermédias ou repetir críticas sem um limite.
\end{promptlabbox}

\section{Quando o sistema escolhe o passo seguinte}

Neste livro, chamamos \emph{agente com LLM} a um sistema em que o modelo pode propor a próxima ação com base num objetivo e no estado observado. O conceito de agente em IA é mais amplo. Aqui, a aplicação decide que ações são admissíveis, aplica permissões e acompanha o resultado. Um modelo pode escolher entre procurar uma política e pedir uma informação, mesmo quando o executor se mantém determinista.

\begin{center}
\small
\begin{tabularx}{\textwidth}{@{}p{.24\textwidth}XX@{}}
\toprule
Desenho & Quem escolhe o passo & Quando experimentar \\
\midrule
Fluxo de regras & O código aplica uma sequência ou um ramo definido. & A ordem é conhecida e pode ser testada. \\
Escolha de rota & Um modelo propõe uma de várias rotas permitidas. & O pedido em linguagem livre exige interpretação. \\
Agente com ferramentas & O modelo propõe a ação seguinte com base no resultado anterior. & Os passos dependem de informação descoberta durante a tarefa. \\
\bottomrule
\end{tabularx}
\end{center}

Um planeador pode propor uma lista de ações. Para a executar, confirma pré-condições e permissões em cada passo. Um supervisor pode distribuir trabalho por especialistas quando as tarefas são separáveis. O capítulo~\ref{ch:multiagente} mostra como medir o custo e verificar as falhas de cada ramo. ReAct estudou a combinação de passos de raciocínio textual com ações em ambientes concretos~\citep{yao2023react}. É um exemplo de arquitetura que podes comparar com uma rota mais simples.

\section{O estado precisa de um dono}

Guarda o objetivo, as fontes usadas, as decisões e as operações que já produziram efeitos. O contexto é o que chega à chamada atual. A memória persistente é informação que a aplicação conserva para uso futuro, com regras de origem, acesso e expiração. Se resumires a conversa, verifica se a síntese conserva a autorização e a última ação confirmada. A secção de engenharia de contexto no capítulo~\ref{ch:context-engineering} trata dessa passagem.

O orquestrador aplica o prazo, o orçamento e as condições de paragem. Distingue trabalho concluído, informação em falta, falha e encaminhamento. Quando um ramo falha, conserva os resultados utilizáveis e marca o que ficou incompleto. A síntese final deve mostrar quais as partes que o sistema conseguiu verificar.

\section{Quando falta a resposta depois de um efeito}

\ficcao Numa extensão hipotética, o serviço cria uma tarefa para um operador. A gravação pode ter acontecido, mas a ligação cai antes de chegar a confirmação. Se o agente voltar a pedir «cria a tarefa» como uma operação nova, o operador pode receber dois pedidos iguais. O cliente conhece o timeout. O estado da tarefa ainda está por confirmar.

Para permitir repetição segura, define uma chave de idempotência para a mesma intenção. No serviço que grava o efeito, regista a chave, os parâmetros relevantes e o resultado de forma consistente com a operação. Uma repetição com a mesma chave e os mesmos parâmetros deve recuperar esse resultado. Se a chave vier com parâmetros diferentes, trata o conflito explicitamente. Quando há dois pedidos em simultâneo, o controlo tem de impedir que ambos criem a tarefa antes de a chave ficar registada.

Se o serviço externo não oferecer esta garantia, tenta consultar o estado por um identificador conhecido ou encaminha a dúvida. Um número de pedido no log ajuda a seguir a chamada. A garantia contra duplicação depende de como o executor trata esse número. Este caso é um desenho de integração. A demonstração da Lavandaria executa consultas e rascunhos, como se pode ver no capítulo~\ref{ch:lavandaria}.
